\scanpage{52}
\margin{\handwritten{\unsure{1|\textbackslash}}}
\mspage{4 - 38 -}
n\textsuperscript{o} 392

où $I$ est pseudo-nul. Il résulte donc du n\textsuperscript{o} 7, prop. 16, (v) et cor. 1 de la prop. 16, que l'on a $c_B(M\otimes_A B)=c_B(Q\otimes_A B)=-c_B(\handwritten{\chi}_B(Q\otimes_A B))=-c_B(i(\handwritten{\chi}_A(Q)))$ en vertu de (ii) ; mais par définition de l'homomorphisme $i:C(A)\longrightarrow C(B)$ on a $c_B(i(\handwritten{\chi}_A(Q)))=i(c_A(\handwritten{\chi}_A(Q)))=-i(c_A(M))$, ce qui achève de prouver (23).

\uline{Remarques}.- 1) Si $M$ est un $A$-module réflexif, $M\otimes_A B$ n'est pas nécessairement réflexif, (exerc. \quad). Il en est toutefois ainsi lorsque $B$ est un $A$-module \uline{plat}, (n\textsuperscript{o} 2, prop. 8).

2) Soit $C$ un troisième anneau noethérien intégralement clos, tel que $B\subset C$ et que $C$ soit un $B$-module de type fini, (donc aussi un $A$-module de type fini). On a alors les formules de transitivité :
\[
N_{C/A}=N_{B/A}\circ N_{C/B}\ ,\tag*{\makebox[0pt][r]{\makebox[\displaywidth][l]{(25)}}}
\]
\[
i_{C/A}=i_{C/B}\circ i_{B/A}\tag*{\makebox[0pt][r]{\makebox[\displaywidth][l]{(26)}}}
\]
qui résultent immédiatement des formules de transitivité pour les indices de ramification et les degrés résiduels, (chap. VI, \S~8, n\textsuperscript{o} 1, lemme 1).

9. \uline{Un théorème de réduction}.

Les notations et hypothèses sont de nouveau celles des n\textsuperscript{os} 2 à 7.

LEMME 6.- \uline{Soient} $R$ \uline{un anneau commutatif}, $\underline p_i$, $(1\leq i\leq n)$, \uline{des idéaux premiers de} $R$, \uline{deux à deux distincts}.

(i) \uline{Pour} $1\leq i\leq n$, \uline{soit} $H_i$ \uline{une partie de} $R/\underline p_i$ \uline{vérifiant la condition suivante} : \uline{il n'existe aucun élément} $\alpha_i\in R/\underline p_i$ \uline{tel que} $\alpha_i+H_i$ \uline{contienne un idéal} $\neq 0$ \uline{de} $R/\underline p_i$. \uline{Alors il existe} $a\in R$ \uline{tel que}, \uline{pour} $1\leq i\leq n$, \uline{l'image canonique de} $a$ \uline{dans} $R/\underline p_i$ \uline{n'appartienne pas à} $H_i$.

(ii) \uline{Si} $\operatorname{Card}(H_i)<\operatorname{Card}(R/\underline p_i)$, \uline{les} $H_i$ \uline{vérifient la condition de} (i).

(i) Raisonnons par récurrence sur $n$, le cas $n=0$ étant trivial. Soit donc $n\geq 1$. Quitte à faire une permutation sur les indices $i$, on peut supposer que $\underline p_1$ est
% archive sheet number: 51
